\section{讨论}
\label{sec:discussion-zh}

共享检查点比较在已实现诊断层面回答了 RQ1。增加候选预算并实施符号评分，在教师强制标记准确率不变的情况下改善了音级集合、音程向量、循环序列顺序、节奏轮廓和姿态指标。因此，这一对比针对完整的引导式解码程序，而不是更强的训练语言模型，也不能被单独归因为重排序。候选选择器不可见的隐藏密度曲线也略微向优选方向变化，但本文没有报告推断不确定性。

条件移除实验更直接地回答了 RQ2。每项移除都在对应控制维度产生了最明确的可解释损失。音级集合和音程向量指标的大幅变化表明，音级标记不仅是固定生成分布的标签；序列、节奏和姿态字段同样提供了其余前缀无法完全恢复的信息。由于语料是程序化的，而且每项比较只使用一个固定种子，这些效应应理解为内部有效性检查，而不是广泛音乐泛化能力的证据。

序列结果揭示了十二音聚合出现与有序实现之间的差别。拟议解码器在序列目标上几乎覆盖完整半音聚合，但精确序列变换准确率为 0，循环序列顺序准确率也只有 0.2411。候选选择能够找到局部或循环一致性更高的片段，却没有解决完整 $P/R/I/RI$ 实现。程序化规则参照接近满分的严格形式指标说明该指标可由直接构造达到，现有差距来自学习生成路径。

RQ3 只在结构记谱层面得到支持。所有尝试导出的 MusicXML 均可解析并匹配请求的小节和声部数，但拟议模型的平均内容跨度比仅为 0.5195。小节数遵循部分来自显式补齐，而不等于持续的音乐活动。因此，这些导出可作为可编辑草图和审计对象，但是否达到出版级制谱以及是否对作曲实践有用，仍需专家评估。

总体而言，结果支持概率生成与显式符号分析之间的分工。神经模型提供变化，确定性理论函数负责揭示并排序约束遵循程度。规则参照在精确序列和节奏实现方面仍更强，说明未来的辅助系统可采用更紧密的混合解码：对部分约束进行精确构造，再让神经生成在剩余自由度内工作。
